%%=============================================================================
%% Vereistenanalyse
%%=============================================================================

\chapter{\IfLanguageName{dutch}{Vereistenanalyse}{Methodology}}%
\label{ch:vereistenanalyse}

\section{Energiebronnen}
Op basis van interviews met onderzoekers van de ET zijn verschillende opties voor energiebronnen geïdentificeerd. Hieronder worden deze verder besproken, met aandacht voor hun toepassing op de ET.

\subsection{SMRs en AMRs}
SMRs en AMRs zijn, volgens de argumentatie in~\ref{sec:energiebronnen_Nucleair}, uitgesloten als potentiële energiebronnen. Deze technologieën kunnen aanzienlijke hoeveelheden energie leveren met een relatief lage ruimtelijke impact. Dit zou ze ideaal maken voor opwekking dicht bij de verbruikers. Het feit dat deze binnen Europa onbestaand zijn en dat de verwachtingen voor de eerste ingebruikname ver na de start van de bouw van de ET ligt maakt dat deze optie niet nader onderzocht wordt. 

\subsection{Windmolens}
De optie voor windmolens is, zoals besproken in~\ref{sec:energiebronnen_Windmolens}, niet haalbaar nabij de geplande locatie. Er is door de verschillende overheden al een perimeter van 10 km rond de geplande locatie uitgezet waarin windmolens niet meer gebouwd mogen worden, er is zelfs al sprake om deze te vergroten naar 15 km. De trillingen veroorzaakt door de windmolens hebben de mogelijkheid om de metingen van de ET te beïnvloeden. Op verdere afstand zijn deze wel mogelijk, opties voor windmolenparken buiten deze perimeter vallen momenteel buiten de scope van het onderzoek. Voor eventuele verdere uitbreiding van de tool zal dit wel in het achterhoofd gehouden moeten worden.

\subsection{Zonne-energie}
Het onderzoeksteam richt zich primair op zonne-energie als een veelbelovend alternatief voor de voorgaande opties. 
Hoewel de energiedichtheid van zonnepanelen per $\mathrm{m}^2$ lager is dan bij andere bronnen \autocite{fritsche2017energy}, bieden zonnepanelen overtuigende voordelen om deze keuze te rechtvaardigen. Zo zijn zonnepanelen flexibel inzetbaar op uiteenlopende locaties, van kleine daken tot grootschalige velden, en kennen ze relatief lage kosten in zowel investering als operationeel beheer \parencite{IREA2025}. Bovendien kunnen ze geïntegreerd worden met bestaand ruimtegebruik, zoals landbouw, wateroppervlakken of daken, zonder extra grond te vereisen. Voor het ET-project is ook het ontbreken van trillingen een cruciaal voordeel, aangezien dit een essentiële vereiste is voor de toepassing.

\section{Datasets}
Door een samenwerking van de ET onderzoeksgroep met HeusschenCopier voor de identificatie van locaties voor de telescoop en het plannen van de proefboringen zijn er al meerdere datasets beschikbaar. Zo bijvoorbeeld de zeer belangrijke dataset waar de geplande hoekpunten van de ET in staan. De websites van de Vlaamse, Belgische, Waalse en Nederlandse overheid bieden ook verschillende kaarten aan. Dit gaat van beschermde natuurgebieden tot landbouwgebieden en hun historische gewassen. De datasets zijn niet gelijk opgebouwd tussen de verschillende instanties wat een uniforme analyse bemoeilijkt.

\section{Benodigde ondersteuning}
De ondersteunde tool die ontwikkeld moet worden zal grafisch de kaarten moeten kunnen weergeven om niet technische stakeholders te kunnen informeren. Ook zal er nood zijn om met dynamische zoekgebieden te kunnen werken, dit zowel om de zoekstraal groter te maken alsook bij veranderlijke zoekgebieden de analyses opnieuw uit te voeren. 

\section{Benodigde toevoeging op tools}
Uit de verschillende gesprekken met de onderzoekers blijkt dat het maken van een alomvattend systeem met beslissingsmechanisme ver buiten de scope van een bachelorthesis valt. De onderzoekers de mogelijkheid geven om gemakkelijk datasets te kunnen analyseren en kwantificeren zal dan ook de voornaamste focus worden. Doordat de onderzoekers een zeer open mandaat hebben om de energiemix te bepalen worden vaak nieuwe opties overwogen. Om te kunnen bepalen of deze opties waarde hebben voor verder onderzoek zal een snelle analyse mogelijk moeten zijn. Hierbij is voornamelijk de beschikbare oppervlakte binnen een zoekgebied nodig. 